\section{重整化与反常}\label{sec:renormalization}
上一节构造了 $\cN{=}2$ 超对称场论的拉格朗日量。在分析它们的动力学之前，先在这里回顾几种基本方法，即单圈重整化与反常。

\subsection{重整化}\index{Renormalization@重整化|hyperpage}\index{OneLoopExactness@单圈精确性|hyperpage}\index{Holomorphy@全纯性|hyperpage}\index{DynamicalScale@动力学能标|hyperpage}\index{BetaFunction@β 函数|hyperpage}\index{ExactlyMarginalCoupling@恰好边缘耦合|hyperpage}\index{DualCoxeterNumber@对偶 Coxeter 数|hyperpage}
回顾一般拉格朗日场论中规范耦合的单圈重整化：
\begin{equation}
\LambdaRG\frac{d}{d\LambdaRG}g=-\frac{g^3}{(4\pi)^2}\left[
\frac{11}{3}C(\text{adj})-\frac23  C(R_f)-\frac13   C(R_s)
\right].
\end{equation}
这里，$\LambdaRG$ 是测量 $g$ 时的能标，并且我们约定将所有费米子都写成左手 Weyl 费米子。于是，$R_f$ 与 $R_s$ 分别为 Weyl 费米子和复标量所属的规范群表示。量 $C(\rho)$ 定义为满足
\begin{equation}
\tr \rho(T^a) \rho(T^b) =C(\rho) \delta^{ab}
\end{equation}
其中 $T^a$ 为规范代数的生成元，$\rho(T^a)$ 是表示 $\rho$ 中的矩阵，归一化使得 $C(\text{adj})$ 等于对偶 Coxeter 数。对于 $\SU(N)$，有
\begin{equation}
C(\text{adj})=N,\qquad C(\text{fund})=\frac12.
\end{equation}
在 $\cN{=}1$ 规范理论中，该方程简化为
\begin{equation}
\LambdaRG\frac{d}{d\LambdaRG}g=-\frac{g^3}{(4\pi)^2} \left[
3C(\text{adj})-C(R)
\right]
\end{equation}
或等价地
\begin{equation}
\LambdaRG\frac{d}{d\LambdaRG}\frac{8\pi^2 }{g^2}=
3C(\text{adj})-C(R),\label{n1running}
\end{equation}
其中 $R$ 是手征多重态的表示。在 $\cN{=}2$ 规范理论中，一个伴随手征多重态 $\Phi$ 被视为矢量多重态的一部分。于是有
\begin{equation}
\LambdaRG\frac{d}{d\LambdaRG}\frac{8\pi^2 }{g^2}=
2C(\text{adj})-C(R),\label{n2running}
\end{equation}
其中 $R$ 现在是描述系统中超多重态的 $\cN{=}1$ 手征多重态的表示。如果存在一个伴随超多重态，由两个 $\cN{=}1$ 手征多重态 $A$ 和 $B$ 组成，那么单圈 $\beta$ 函数为零。当 $A$、$B$ 的质量项为零时，系统实际上具有进一步扩大的超对称性，成为 $\cN{=}4$ 超杨--米尔斯理论。当质量项非零时，则称为 $\cN{=}2^*$ 理论。

在超对称理论中，耦合 $g$ 与 $\theta$ 角 $\theta$ 合并，并以下述形式进入拉格朗日量：
\begin{equation}
\ \int d^2\theta \frac{-i}{8\pi}\tau \tr W_\alpha W^\alpha + cc.
\end{equation}
其中 $\tau$ 为
\begin{equation}
\tau=\frac{4\pi i}{g^2} +\frac{\theta}{2\pi}.
\end{equation}
我们将 $\tau$ 称为复化规范耦合。

可以把 $\tau$ 看作某个背景手征超场的期望值。存在一种重整化方案，使超势始终是手征超场的全纯函数；这里也包括真空期望值给出规范耦合和超势耦合的背景场\cite{Intriligator:1995au}。我们称其为 Seiberg 全纯性原理。

在这一方案中，能标 $\LambdaRG$ 处的单圈跑动耦合可以表示为
\begin{equation}
\tau(\LambdaRG )=\tau_{UV} - \frac{b}{2\pi i} \log\frac{\LambdaRG }{\Lambda_{UV}} + \cdots
\label{1looprunning}
\end{equation}
其中 $b$ 是式\eqref{n1running}或式\eqref{n2running}右边出现的有理数。注意，耦合 $\tau$ 从 $1/g^2$ 开始，因此 $n$ 圈图会具有 $g^{2(n-1)}$ 的依赖关系。于是，式\eqref{1looprunning}中虚部的常数平移就是一个单圈效应。

微扰论与 $\theta$ 角无关，因为 $F_{\mu\nu}\tilde F_{\mu\nu}$ 是全导数，尽管它是一个规范依赖量的全导数。因此，$n$ 圈效应是 $(\Im\tau)^{1-n}$ 的函数，除非 $n=1$，否则它不全纯。我们得出结论，在全纯方案中，跑动关系\eqref{1looprunning}是单圈精确的。可以发现，组合
\begin{equation}
\Lambda^b = \LambdaRG ^b e^{2\pi i  \tau(\LambdaRG )}
\end{equation}
在微扰论的所有阶都不变。我们将 $\Lambda$ 称为理论的复化动力学能标。\footnote{以实常数 $c$ 作形式为 $\Lambda\to c\Lambda$ 的重新定义，对应于形式为 $1/g^2\to1/g^2-c'$ 的耦合重新定义，其中 $c'$ 为另一常数；或等价地，$g^2\to g^2+c'g^4+\cdots$。因此，这一重新定义始于单圈阶，并保持 $g^2$ 的领头阶定义不变。本讲义并不十分仔细地追踪这类耦合的有限重整化。} 注意，$\Lambda$ 是复量，可以看作背景手征超场的真空期望值。

这种单圈精确性并不必然意味着，例如控制散射过程的物理规范耦合也是单圈精确的。对于一般的 $\cN{=}1$ 超对称理论，在全纯方案中存在非平凡的波函数重整化因子 $Z_{ij}$：
\begin{equation}
\int d^4\theta Z^{\bar ij}(\LambdaRG ) Q^\dagger_{\bar i}  e^{V} Q_j \label{wr}
\end{equation}
计算物理散射振幅时，需要通过进一步的场重新定义将这些因子计入。已知这会对规范耦合的物理跑动产生额外的微扰贡献。更多讨论例如可参见文献\cite{ArkaniHamed:1997mj}。

不过，对于 $\cN{=}2$ 超对称理论，可以作出更强的论断。我们假设存在一种全纯重整化方案，它与 $\SU(2)_R$ 对称性的存在相容。此时，拉格朗日量的结构受到限制：矢量多重态部分必须具有式\eqref{vectorlag}的形式，超多重态部分必须具有式\eqref{hyperlag}的形式。将 $\tau$ 视为背景场的真空期望值。于是，在矢量多重态一侧，可以发现拉格朗日量\eqref{vectorlag}中不能出现式\eqref{wr}那样非平凡的波函数重整化因子 $Z_{\bar ij}$。在超多重态一侧，式\eqref{hyperlag}中的系数 $c'$ 在全纯方案中不被重整化。由于 $c=c'$，凯勒势也不被重整化。因而，超多重态拉格朗日量\eqref{hyperlag}中不存在重整化。

于是，特别是在 $b=0$ 时，$\beta$ 函数在微扰论所有阶都为零。这使系统具有共形性，而 $\tau$ 的值成为一个恰好边缘的耦合参数。非微扰修正会引起有限重整化，但人们认为它们不会引入任何额外的无限重整化。

例如，$\cN{=}4$ 超杨--米尔斯理论自动具有超共形性，并有一个恰好边缘耦合。另一个 $b=0$ 的例子，是含有 $2N$ 个基本表示超多重态的 $\cN{=}2$ 超对称 $\SU(N)$ 规范理论。事实上，在式\eqref{n2running}中，有 $C(\text{adj})=N$，$C(R)=2\cdot2N\cdot1/2$。

\subsection{反常}
\subsubsection{整体对称性的反常}\label{sec:anomaly}\index{Anomaly@反常|hyperpage}\index{Instanton@瞬子|hyperpage}\index{InstantonNumber@瞬子数|hyperpage}
非阿贝尔规范理论中的非微扰效应有一个重要来源，称为瞬子。它是欧氏 $\mathbb{R}^4$ 上的非平凡经典场组态，具有非零积分
\begin{equation}
16\pi^2 k :=\int_{\mathbb{R}^4} \tr F_{\mu\nu}\tilde F^{\mu\nu}.
\end{equation}
在 $\SU(N)$ 的迹的标准归一化下，$k$ 自动为整数，称为瞬子数。欧氏路径积分中的 theta 项以如下形式出现：
\begin{equation}
\exp\left[i\frac{\theta}{16\pi^2} \tr F_{\mu\nu} \tilde F^{\mu\nu}\right].
\end{equation}
因此，瞬子数为 $k$ 的组态具有非平凡相位 $e^{i\theta k}$。注意，将 $\theta$ 平移 $2\pi$ 完全不会改变这一相位。因而，即使在量子理论中，平移 $\theta\to\theta+2\pi$ 也仍然是一种对称性。

利用
\begin{equation}
\tr F_{\mu\nu}F_{\mu\nu} = \frac12 \tr(F_{\mu\nu}\pm \tilde F_{\mu\nu})^2 \mp  \tr F_{\mu\nu} \tilde F_{\mu\nu}
\ge \mp \tr F_{\mu\nu}\tilde F_{\mu\nu},
\end{equation}
得到
\begin{equation}
\int d^4 x \tr F_{\mu\nu} F_{\mu\nu} \ge 16\pi^2 |k| \label{bound}
\end{equation}
它只在
\begin{equation}
F_{\mu\nu} + \tilde F_{\mu\nu}\propto F_{\alpha\beta}=0 \quad
\text{或}
\quad F_{\mu\nu} - \tilde F_{\mu\nu}\propto F_{\dot\alpha\dot\beta}=0
\label{SD}
\end{equation}
时取等号，具体取决于 $k$ 的符号。因此，在固定 $k$ 的组态中，满足关系\eqref{SD}的组态对路径积分给出主导贡献。根据 $k$ 的符号，式\eqref{SD}的解称为瞬子或反瞬子。

在瞬子背景中，路径积分中来自规范动能项的权重为
\begin{equation}
\exp\left[
-\frac{1}{2g^2}\int \tr F_{\mu\nu}F^{\mu\nu}
+
i\frac{\theta}{16\pi^2}\int \tr F_{\mu\nu}\tilde F^{\mu\nu}
\right]
=e^{2\pi i \tau k}.\label{instanton_contri}
\end{equation}
类似地，在反瞬子背景中有贡献 $e^{2\pi i\bar\tau |k|}$。这里仅出现 $\tau$ 或 $\bar\tau$，而不是更复杂的组合，这与下述事实有关：在超对称理论的瞬子背景中，假定 D 项也为零，则有 $\delta\lambda_{\dot\alpha}=F_{\dot\alpha\dot\beta}\epsilon^{\dot\beta}=0$，因此带点的超平移被保持。类似地，在反瞬子背景中，不带点的超对称性保持不破缺。

现在，考虑与规范场耦合的带荷 Weyl 费米子 $\psi_\alpha$，其动能项为
\begin{equation}
\bar\psi_{\dot\alpha} {D}{}_\mu \sigma^{\mu \dot\alpha\alpha} \psi_\alpha.
\end{equation}
设 $\psi_\alpha$ 属于规范群的表示 $R$。已知 $\psi_\alpha$ 中的零模数减去 $\bar\psi_{\dot\alpha}$ 中的零模数，等于 $2C(R)k$。特别地，在 $k$ 为正时，限制到 $k$ 瞬子组态的路径积分为零，除非在被积函数中额外插入 $k$ 个 $\psi$。更明确地说，
\begin{equation}
\vev{O_1O_2\cdots} = \int [D\psi] [D\bar\psi] O_1O_2 \cdots e^{-S} =0
\end{equation}
除非算符乘积 $O_1O_2\cdots$ 中包含的 $\psi$ 比 $\bar\psi$ 多 $2C(R)k$ 个。这可以解释如下：路径积分测度 $[D\psi]$ 和 $[D\bar\psi]$ 都包含无穷多个积分。不过，两者的积分变量数目相差一个有限数 $2C(R)k$。等价地，在常数转动
\begin{equation}
\psi \to e^{i\varphi} \psi,\quad
\bar\psi \to e^{-i\varphi} \bar\psi,
\end{equation}
下，费米子路径积分测度按照
\begin{equation}
\begin{aligned}\relax
[D\psi] & \to [D\psi]e^{+\infty i\varphi + 2C(R)k i\varphi  }, \\
[D\bar\psi] & \to [D\bar\psi]e^{-\infty i\varphi }.
\end{aligned}
\end{equation}
转动。将二者合并，得到
\begin{equation}
[D\psi][D\bar\psi] \to [D\psi][D\bar\psi] e^{2C(R)k i\varphi} =
[D\psi][D\bar\psi] \exp\left[
2C(R)\varphi \frac{i}{16\pi^2}\int \tr F_{\mu\nu}\tilde F^{\mu\nu}
\right].
\end{equation}
这可以通过平移 $\theta$ 角，即 $\theta\to\theta+2C(R)\varphi$ 来补偿。正如前面所回顾的，平移 $\theta\to\theta+2\pi$ 是一种对称性。因而，对场 $\psi$ 施以 $\exp(\frac{2\pi i}{2C(R)})$ 的转动，是一种真正未破缺的对称性。

\subsubsection{规范对称性的反常}\index{GaugeAnomaly@规范反常|hyperpage}\index{WittenGlobalAnomaly@Witten 整体反常|hyperpage}
在 $\cN{=}2$ 规范理论中，费米子总是出现在非手征表示中。事实上，矢量多重态中的费米子总属于伴随表示，完整超多重态中的 $\cN{=}1$ 手征超场构成表示 $R$ 与其共轭表示 $\bar R$ 之和，而一个半超多重态则相当于赝实表示 $R$ 中的一个 $\cN{=}1$ 手征超场。因此，不存在微扰规范反常。

不过，必须小心 Witten 整体反常\cite{Witten:1982fp}，因为即使对于实表示，它也可能出现。已知规范群 $\SU(2)$ 二重态中的一个 Weyl 费米子具有反常，其原因如下。进行这一系统的路径积分时，首先需要考虑
\begin{equation}
Z[A_\mu] = \int [D\psi_{\alpha i}][D\bar\psi_{\dot\alpha i}] e^{-\int \bar\psi D_\mu \sigma^\mu \psi}
\end{equation}
其中 $i=1,2$ 是 $\SU(2)$ 二重态指标。为了进一步一致地对 $A_\mu$ 积分，需要
\begin{equation}
Z[A_\mu]= Z[A^g_\mu], \quad A^g_\mu = g^{-1} A_\mu g  + g^{-1} \partial_\mu g.
\end{equation}
对任意规范变换 $g:\mathbb{R}^4\to\SU(2)$ 都成立。这些映射由 $\pi_4(\SU(2))$ 刻画。已知
\begin{equation}
\pi_4(\SU(2))=\pi_4(S^3)=\mathbb{Z}_2.
\end{equation}
令 $g_0:\mathbb{R}^4\to\SU(2)$ 为对应于这一 $\bZ_2$ 中非平凡元素的映射。则已知
\begin{equation}
[D\psi_{\alpha i}][D\bar\psi_{\dot\alpha i}] \stackrel{g_0}{\longrightarrow}
-[D\psi_{\alpha i}][D\bar\psi_{\dot\alpha i}]
\end{equation}
从而有
\begin{equation}
Z[A^{g_0}_\mu] = -Z[A_\mu],
\end{equation}
使得对 $A_\mu$ 的路径积分不一致。

一般而言，若 $G=\Sp(n)$，则 $\pi_4(G)=\mathbb{Z}_2$；否则 $\pi_4(G)=1$。因此，只有规范群 $\Sp(n)$ 的表示 $R$ 中的 Weyl 费米子才可能存在 Witten 整体反常。一个简短的计算表明，只有当 $C(R)$ 为半整数时，才存在反常。

在四维中，Witten 反常总是取值于 $\bZ_2$。因而，完整超多重态总是不具有 Witten 整体反常。这种危险只存在于规范群 $\Sp(n)$ 的半超多重态中。例如，不能有奇数个规范群 $\SU(2)$ 二重态表示中的半超多重态；更一般地，不能有规范群 $\Sp(n)$ 的赝实表示 $R$ 中满足 $C(R)$ 为半整数的半超多重态。

\subsection{$\cN{=}1$ 纯杨--米尔斯理论}\label{sec:n=1pure}
\subsubsection{禁闭与规范微子凝聚}\index{Confinement@禁闭|hyperpage}\index{GauginoCondensate@规范微子凝聚|hyperpage}
作为本节所学内容的一个应用，考虑规范群为 $\SU(N)$ 的 $\cN{=}1$ 纯超对称杨--米尔斯理论。这一小节的内容在本讲义后面不会被大量使用。

该理论只有矢量多重态，拉格朗日量为
\begin{equation}
L=\int d^2\theta \frac{-i}{8\pi}\tau \tr W_\alpha W^\alpha+ cc., \qquad
W_\alpha=\lambda_\alpha + F_{\alpha\beta}\theta^\beta + \cdots
\end{equation}
耦合的单圈跑动为
\begin{equation}
E\frac{\partial}{\partial E}\tau(E) = 3N,
\end{equation}
% [Tutor Note: p.31 式 (3.3.2) 原文（v2）作 $E\frac{\partial}{\partial E}\tau(E)=3N$，缺因子 $\frac{i}{2\pi}$：
%  由式 (1looprunning) 的 $\tau(E)=\tau_{UV}-\frac{b}{2\pi i}\log\frac{E}{\Lambda_{UV}}$ 且纯 $\cN{=}1$ SYM 的 $b=3N$，
%  得 $E\frac{\partial}{\partial E}\tau(E)=\frac{i}{2\pi}3N$；也可由 $\Lambda^{3N}=E^{3N}e^{2\pi i\tau(E)}$ 对 $\log E$ 求导
%  （$3N+2\pi i\,E\tau'=0$）得到同一结果。arXiv v3 源码已把此处改为 $\frac{i}{2\pi}3N$。正文照 v2 原文排。]
因而通过关系
\begin{equation}
\Lambda^{3N}= e^{2\pi i\tau_{UV} }\Lambda_{UV}^{3N}.
\end{equation}
定义动力学能标 $\Lambda$。

赋予规范场零 R 荷，赋予规范微子 $\lambda_\alpha$ 的 R 荷为 1。相位转动 $\lambda_\alpha\to e^{i\varphi}\lambda_\alpha$ 存在反常，需要用 $\theta\to\theta+2N\varphi$ 来补偿。将 $\theta$ 平移 $2\pi$ 仍是一种对称性，因此离散转动
\begin{equation}
\lambda_\alpha \to e^{\pi i/N} \lambda_\alpha, \qquad \theta\to \theta+2\pi
\end{equation}
是一种生成 $\bZ_{2N}$ 的对称性。注意，在这一对称性下，上面定义的 $\Lambda$ 按照
\begin{equation}
\Lambda \to e^{2\pi i/(3N)} \Lambda.\label{lambdatrans}
\end{equation}
变换。

人们认为，这一理论发生禁闭，并具有非零的规范微子凝聚 $\vev{\lambda_\alpha\lambda^\alpha}$。这一凝聚的值会是多少？它应当具有质量量纲 3 和 R 荷 2。唯一的候选是
\begin{equation}
\vev{\lambda_\alpha\lambda^\alpha}=c \Lambda^3
\end{equation}
其中 $c$ 为某个常数。对称性\eqref{lambdatrans}在两边都以乘上 $e^{2\pi i/N}$ 的相同方式作用。假定数值常数 $c$ 非零，$\bZ_{2N}$ 会进一步自发破缺为 $\bZ_2$，产生 $N$ 个不同的解：
\begin{equation}
\vev{\lambda_\alpha\lambda^\alpha}=c e^{2\pi i \ell/N } \Lambda^3
\end{equation}
其中 $\ell=0,1,\ldots,N-1$。未破缺的 $\bZ_2$ 以 $\lambda_\alpha\to-\lambda_\alpha$ 作用于费米子，这对应于 $360^\circ$ 转动。这一 $\bZ_2$ 对称性很难破缺。

现在普遍认为，这一理论恰好具有上述 $N$ 个超对称真空，没有更多真空。对其他规范群，分析方式相同，只需用所考虑规范群的对偶 Coxeter 数 $C(\text{adj})$ 替换 $N$。例如，纯 $\cN{=}1$ $\SO(N)$ 规范理论具有 $N-2$ 个真空。

\subsubsection{盒子中的理论}\label{sec:box}\index{WittenIndex@Witten 指标|hyperpage}\index{Holonomy@和乐|hyperpage}
回顾另一种计算规范群为 $G$ 的 $\cN{=}1$ 纯杨--米尔斯理论真空数目的方法，是很有启发性的；这一方法最初在文献\cite{Witten:1982df}中讨论。将系统放入一个大小为 $L\times L\times L$ 的空间盒子，并在每个方向施加周期性边界条件。时间方向仍保持为 $\bR$。沿三个空间方向作 Kaluza-Klein 约化，系统就成为具有无穷多个自由度的超对称量子力学。

这一盒子仍然保持平移生成元 $P^\mu$ 和超平移 $Q_\alpha$ 不破缺。我们只使用 $Q_\alpha$ 与 $Q_\alpha^\dagger$ 的一个线性组合 $\cQ$，满足
\begin{equation}
H=P^0=  \{\cQ,\cQ^\dagger\}.
\end{equation}
还存在费米子数算符 $(-1)^F$，使得
\begin{equation}
\{(-1)^F,\cQ\}=0.
\end{equation}
考虑哈密顿量 $H$ 的本征态，满足
\begin{equation}
H\ket{E}=E\ket{E}.
\end{equation}
一般而言，在 $\cQ$、$\cQ^\dagger$、$H$ 与 $(-1)^F$ 的代数下，多重态结构具有如下形式：
\begin{equation}
\begin{array}{ccccc@{\qquad}c}
&\raisebox{-.8ex}{\rotatebox{30}{$\leftrightarrow$}}& \cQ^\dagger \ket{E} & \leftrightarrow & (\cQ^\dagger\cQ-\cQ\cQ^\dagger)\ket{E} \\
\ket{E} & \leftrightarrow & \phantom{^\dagger}\cQ\ket{E} &  \rotatebox{30}{$\leftrightarrow$}
\end{array}
\end{equation}
其中包含四个态。当 $\cQ\ket{E}=0$ 或 $\cQ^\dagger\ket{E}=0$ 时，多重态只有两个态。若 $\cQ\ket{E}=\cQ^\dagger\ket{E}=0$，多重态便只有一个态，而且由等式
\begin{equation}
E\vev{EE}
=\vev{E|H|E}
=\vev{E|(\cQ \cQ^\dagger+\cQ^\dagger \cQ)|E}=
|\cQ\ket{E}|^2+
|\cQ^\dagger\ket{E}|^2 .
\end{equation}
可知 $E$ 自动为零。可以看到，除非 $E=0$，否则一个玻色态总会与一个费米态配对。

这保证了 Witten 指标
\begin{equation}
\tr e^{-\beta H} (-1)^F =\tr\big|_{E=0} (-1)^F
\end{equation}
是一个不随盒子大小 $L$ 改变的稳健量：当某个扰动使若干零能态变为非零能量 $E\neq0$ 的态时，所涉及的态必定由成对的费米态和玻色态构成。因而，它不能改变 $\tr(-1)^F$。

因此，可以在盒子大小 $L$ 远小于动力学所确定的尺度 $\Lambda^{-1}$ 的极限下计算 Witten 指标。此时系统弱耦合，可以使用微扰分析。为了使能量接近零，需要空间方向上的 $F_{\mu\nu}=0$，因为磁场会对能量作贡献。于是，系统中唯一的低能自由度就是彼此对易的和乐
\begin{equation}
U_x, U_y, U_z \in \SU(N),
\end{equation}
假设它们可以同时对角化，则有
\begin{align}
U_x&=\diag(e^{i\theta^x_1},\ldots,e^{i\theta^x_N}),\\
U_y&=\diag(e^{i\theta^y_1},\ldots,e^{i\theta^y_N}),\\
U_z&=\diag(e^{i\theta^z_1},\ldots,e^{i\theta^z_N}).\label{Uxyz}
\end{align}
同时还有规范微子零模
\begin{equation}
\lambda^{\alpha=1}_1,\ldots,\lambda^{\alpha=1}_N,\quad
\lambda^{\alpha=2}_1,\ldots,\lambda^{\alpha=2}_N
\end{equation}
并满足条件
\begin{equation}
\sum_i \theta^x_i=\sum_i \theta^y_i=\sum_i \theta^z_i=0,\quad
\sum_i\lambda^{\alpha=1}_i=
\sum_i\lambda^{\alpha=2}_i=0.
\end{equation}
这一截断后的量子系统的波函数，是具有如下形式的态的线性组合：
\begin{equation}
\lambda^{\alpha_1}_{i_1}\lambda^{\alpha_2}_{i_2}\cdots \lambda^{\alpha_\ell}_{i_\ell}\psi(\theta^x_i;\theta^y_i;\theta^z_i)
\end{equation}
它在作用于指标 $i=1,\ldots N$ 的置换下不变。为了具有零能量，波函数不能依赖 $\theta^{x,y,z}_i$，因为对这些变量的导数就是电场分量，会对能量作贡献。因而，唯一可能的零能态就是 $\lambda$ 的不变多项式。我们找到 $N$ 个态，其波函数为
\begin{equation}
1,\ S,\ S^2,\ \ldots,\ S^{N-1}\label{tow}
\end{equation}
其中 $S=\sum_i\lambda^{\alpha=1}_i\lambda^{\alpha=2}_i$。它们都具有相同的 Grassmann 奇偶性，并以相同符号对 Witten 指标作贡献。因此，在小盒子极限中也找到了 $N$ 个态。

将上述构造应用到其他群，只会得到 $1+\rank G$ 个态。例如，考虑 $N>4$ 时的 $G=\SO(N)$。目前介绍的方法只给出 $\lfloor N/2\rfloor+1$ 个态：
\begin{equation}
1,\ S,\ S^2,\ \ldots,\ S^{\lfloor N/2\rfloor},
\end{equation}
而在 $N\ge7$ 时，这与 $C(\text{adj})=N-2$ 不一致。这一难题已经在文献\cite{Witten:1982df}中指出，后来由同一作者在文献\cite{Witten:1997bs}的附录 I 中解决。\footnote{文献\cite{Witten:1982df}如此重要，其中提出的问题却在 15 年间都未被其他物理学家解决，这实在令人遗憾。看来，我们这一领域的人们过于依赖文献\cite{Witten:1982df,Witten:1997bs}的作者了。} 错误之处在于，假设三个彼此对易的矩阵 $U_{x,y,z}$ 可以像式\eqref{Uxyz}那样同时对角化。已知还存在另一个分支，在其中它们不能同时对角化到 Cartan 环面中。对于 $\SO(7)$，下面这个三元组就是一个例子：
\begin{align}
U_x^{(7)}&=\diag(++---+-),\\
U_y^{(7)}&=\diag(+-+-+--),\\
U_z^{(7)}&=\diag(-+++---).
\end{align}
这三个矩阵看起来可能是对角的，但并不属于同一个 Cartan 子群。这一分支增加一个超对称态。因此，总共有 $(\lfloor7/2\rfloor+1)+1=5=7-2$ 个态，再现了 $C(\text{adj})$。

对于更大的 $N$，可以考虑如下形式的 $U_{x,y,z}$：
\begin{equation}
U_x=U_x^{(7)} \oplus U_x',\quad
U_y=U_y^{(7)} \oplus U_y',\quad
U_z=U_z^{(7)} \oplus U_z',\label{N-7}
\end{equation}
其中 $U_{x,y,z}'$ 属于 $\SO(N-7)$ 的 Cartan 子群。在两个分支中分别应用导出式\eqref{tow}的分析，即分别考虑 $U_{x,y,z}$ 属于 $\SO(N)$ 的 Cartan 子群的分支，以及 $U_{x,y,z}$ 具有式\eqref{N-7}形式的分支，总共得到
\begin{equation}
(\lfloor N/2\rfloor+1)+(\lfloor (N-7)/2\rfloor+1) = N-2
\end{equation}
个零能态，从而再现了 $C(\text{adj})$ 个态。这一分析已被推广到任意规范群\cite{Kac:1999gw,BorelFreedmanMorgan}。
